\chapter{Einordnung: Wo die Horimetrik wohnt}
\label{ch:einordnung}

% Register: S44, V13. Quellen: Anhang J, experiments_einordnung.py.
% Schlusskapitel: erst die Fronten (K8), dann die Landkarte.

Ein Buch über eine neue Theorie schuldet seinen Lesern am Ende eine
Adresse: Wo auf der Landkarte der Mathematik liegt das alles? Die
Antwort hat sich erst ganz zum Schluss der Arbeit vollständig
zusammengesetzt, und sie ist besser, als wir zu hoffen wagten. Es
gibt nämlich drei mögliche Ausgänge für ein Projekt wie dieses. Ein
\emph{unbekannter} Grundraum wäre verdächtig -- vermutlich hätte man
nur schlecht gesucht. Eine \emph{bekannte vollständige Theorie} wäre
das Ende -- man hätte alte Mathematik neu angemalt. Der beste Fall
ist der dritte: ein bekannter Grundraum, auf dem neue Operationen
und neue Sätze stehen. Genau dort sind wir gelandet.

\section{Die Bewohner sind alte Bekannte}

Ergänzt man die Ziffern einer Hori-Zahl um eine führende Null,
$e=(0,a_1,\ldots,a_n)$, so gilt an jeder Stelle $0\le e_j<j$ -- und
das ist Wort für Wort die Definition einer
\emph{Inversionsfolge} der Länge $n{+}1$. Inversionsfolgen stehen
in klassischer Bijektion zu den Permutationen: $e_j$ zählt, wie
viele frühere Einträge einer Permutation größer sind als der
$j$-te. Also:

\begin{satz}[Register S44]
$H_n$ steht in Bijektion zu den Inversionsfolgen der Länge $n{+}1$
und damit zu den Permutationen von $n{+}1$ Elementen. Unter dieser
Brille wird der Strukturhorizont zur Permutationsstatistik
\[
\sigma(P_a)=n-\min_{e_j>0}\bigl((j{-}1)-e_j\bigr),
\]
und seine Verteilung ist exakt bekannt: Auf jedem Niveau
$1\le r\le n$ liegen genau $r\cdot r!$ Elemente (die Formel gilt
für $a\ne0$; die Null hat $\sigma=0$).
\end{satz}

Die Fakultätsgröße der Horizonte war also nie ein Zufall: Die
Bewohner der Hori-Welten \emph{sind} Permutationscodes. Das nimmt
der Theorie nichts -- es sagt präzise, was sie ist: Die Horimetrik
untersucht, wie diese wohlbekannten Codes ihren \emph{Wert}, ihre
\emph{Gestalt} und ihre \emph{Identität} verändern, wenn ihr Raum
wächst. Umzüge, Projektoren, Seitentreue, Kapazität -- all das ist
auf dem alten Grundraum neu.

Die Statistik selbst hat eine hübsche Lesart: $(j{-}1)-e_j$ ist der
Abstand einer Ziffer von ihrer oberen Schranke. Die Literatur kennt
die Extremfälle als \emph{tight entries} (Ziffern, die ganz oben
anliegen); der Strukturhorizont verfeinert das zu einem Maß: Wie
nahe kommt \emph{irgendeine} Ziffer ihrem Anschlag? Ob diese
Statistik unter der Bijektion einer klassischen
Permutationsstatistik entspricht, ist offen (Register V13) -- ihre
Verteilung $r\cdot r!$ ist jedenfalls fast verdächtig sauber.

\section{Die Ordnung: ein Gitter, das es schon gibt}

Vergleicht man Ziffernfolgen stellenweise, wird jeder Horizont zu
einem Produkt endlicher Ketten -- einem \emph{distributiven
Gitter}. In dieser Ordnung ist die randvolle Gestalt $M_r$ ein
gewöhnliches Element, und die Kapazitätsklasse
$\mathcal F_r=\{P:\sigma(P)\le r\}$ ist schlicht das
\emph{Hauptideal} darunter: alles, was unter $M_r$ liegt. Das ist
die sauberste abstrakte Beschreibung des Strukturhorizonts, die wir
kennen: $\sigma(P)$ ist die erste Stufe der Kette
$M_0\le M_1\le M_2\le\cdots$ (mit $M_0=0$), deren Hauptideal $P$
enthält -- und
das Extremalprinzip aus Kapitel~\ref{ch:strukturpolynome} wird in
dieser Sprache beinahe eine Tautologie. Dieselbe stellenweise
Ordnung auf Inversionsfolgen wurde jüngst als \emph{middle order}
auf Permutationen untersucht, zwischen schwacher und
Bruhat-Ordnung. Ob unsere Projektoren $K_V$, $K_S$ in diesem Gitter
monoton sind, ob der Kommutator ein Ordnungsdefekt ist -- das sind
ausformulierte Anschlussfragen an ein existierendes Gebiet.

\section{Der Kalkül: alte Operatoren, neue Buchhaltung}

Die fallenden Fakultäten sind die Hausbasis der
\emph{finiten Operatorrechnung}: $\Delta$ spielt dort die Rolle der
Ableitung, und die umbrale Algebra studiert seit Jahrzehnten
Delta-Operatoren und ihre Basisfolgen. Strukturpolynome und
Horizontableitung sind also bestens verortet. Was die Horimetrik
hinzufügt, ist eine Frage, die dort nicht im Mittelpunkt steht:
nicht \emph{wie} ein Operator ein Polynom verändert, sondern
\emph{wie viel minimale Strukturkapazität} das kostet. Das
Horizontkalkül -- $\Gamma_+$, $\Gamma_\times$, $\Gamma_\Delta$,
$\Gamma_E$, $\Gamma_\circ$, alle exakt und alle am selben
kanonischen Objekt angenommen -- ist eine Kapazitätsbuchhaltung
über einem etablierten Operatorgerüst. Zusammen mit der
Zählbreiten-Deutung aus Kapitel~\ref{ch:strukturpolynome} ist das
der Kandidat für einen eigenständigen Fachaufsatz ganz ohne
Hori-Erzählung.

\section{Das Forschungsprogramm dahinter}

Drei weitere Nachbarschaften haben wir gesichtet, aber nicht mehr
betreten; sie seien als Programm notiert. \emph{Erstens:} Ersetzt
man die Ziffernschranken $0\le a_i\le i$ durch beliebige
Kapazitätsprofile $0\le a_i<s_i$, entstehen die in der Kombinatorik
untersuchten $s$-Inversionsfolgen samt ihrer Geometrie
(Lecture-Hall-Polytope). Eine \emph{$s$-Horimetrik} wäre eine
zweite Theorie -- welche Profile eine Polynombasis, eine
Wert-Gestalt-Dualität, eine Seitentreue erlauben, ist vollständig
offen. \emph{Zweitens:} Die Rook-Theorie schreibt ihre Polynome in
genau unserer Basis; welche Gestalten Rook-Polynome von Brettern
sind und ob $M_r$ ein natürliches Brett besitzt, wäre ein
kombinatorischer Zugang zu unseren Produktformeln.
\emph{Drittens:} Die Zählsemantik der Gestalten -- gefärbte
geordnete Injektionen -- ist der Rohstoff der kombinatorischen
Spezies; dort könnte eine gemeinsame Sprache für alle
Interpretationen dieses Buchs liegen.

\begin{oberflaeche}
Das Fazit des Buchs, in drei Zeilen: Die Bewohner der Horimetrik
sind bekannte Objekte -- Permutationscodes in einem bekannten
Gitter, beschrieben in der Hausbasis der Differenzenrechnung. Die
Operationen darauf -- zwei Arten zu wachsen, zwei Arten
aufzuräumen, eine Kapazitätsbuchhaltung, eine Zählbreite -- sind
neu, und die Sätze über sie auch. Ein bekannter Grundraum mit
neuen Operationen und neuen Sätzen: Für eine kleine Theorie, die
mit dem Umdrehen einer Ziffernordnung im Dunkeln begann, ist das
vermutlich der beste Platz, den die Landkarte zu vergeben hat.
\end{oberflaeche}
